\section{Будущая работа}
\label{sec:future}

Будущая работа \DEIC{} распределена по двум осям: технические расширения (методологические, не требующие теоретической переработки) и фундаментальные продолжения (разработка теоретических гапов из §\ref{sec:limitations}, каждый из которых выводится в отдельную статью программы).

\subsection{Технические расширения}

\paragraph{Pre-registered replication на новой выборке.} Обязательное условие до wide publication. Центральные предсказания --- A-модерация канала CA$\to$ID$\to$\EAF{} в узкой спецификации, коллапс ID-пути в unified SEM при включении P, двухосевое различение первичного и вторичного паттернов, компенсаторная функция A в hi-CA$\times$hi-A группе, положительная partial-корреляция P и \Ecog{} после partialling. Replication должна воспроизвести именно эту комбинацию, включая ре-композицию дисперсии в unified модели.

\paragraph{Item wave 2: инструментальная пересборка.} Приоритетные направления: (1) dedicated A-scale, разведённая по дизайну между witnessing и trauma hypervigilance, с контентной валидизацией на cover всех четырёх лингвистических традиций (satipaṭṭhāna, sākṣī, rigpa, self-remembering); (2) M-редизайн с разведением instrumental и reactive masking; (3) \EAF{}/\Ecog{} разведены на item-level, а не post-hoc; (4) один айтем $=$ один фактор в scoring (без multi-scale весов); (5) P-шкала расширена явной моральной подписью (items о «разрешено мне») отдельно от архитектурных items о resonance switching; (6) ACE-айтемы сохраняют трёхвариантный формат («да» / «нет» / «не помню»), но «не помню» кодируется как эксплицитный \emph{позитивный} сигнал (семейное замалчивание или диссоциативный encoding-gap), а не missing; счётчик таких ответов по всем 10 айтемам формирует pre-registered ACE-dissociation подшкалу с заявленными критериями McDonald's $\omega \ge 0.7$ и $r(\text{ACE-dis}, \text{ID}_{\mathrm{lat}}) \ge 0.30$. Post-hoc re-fit на wave-1 (§\ref{subsec:ace_extension}) показывает усиление латентной связи $r(\text{CA}, \EAF{})$ с $-0.129$ до $-0.205$ и появление статистически отличимого от нуля индиректа через P --- wave-2 должна pre-register эту логику и подтвердить её как установленный факт.

\paragraph{Лонгитюдное расширение} (6--12--24 месяца, in-therapy subgroup). Тестирует стабильность A, траекторию integration, test-retest-надёжность всех шести факторов. Выборка с явной стратификацией по therapy-status (in-treatment, post-treatment, never-treated) позволит различить стабильное и транзиторное A.

\paragraph{Мультикультурная валидация} (EN, ES, CH, AR) с invariance-testing (configural $\to$ metric $\to$ scalar). Проверка обобщаемости архитектуры за пределами русскоязычной выборки.

\paragraph{fMRI-конвергенция.} DMN $\times$ attention-networks co-activation как нейро-коррелат A. Биологическая грануляция A-конструкта через существующие neuroimaging-парадигмы (resting-state functional connectivity, mindfulness-induced DMN-deactivation).

\paragraph{Informant report на субвыборке.} Круговая валидность особенно важна для A и P, где self-report имеет структурные ограничения.

\paragraph{Behavioral convergents для A.} Dual-task paradigms, interoception accuracy tasks (heartbeat detection), mind-wandering probes --- выход из чистого self-report.

\subsection{Фундаментальные продолжения}

Это статьи программы, разрабатывающие теоретические гапы §\ref{sec:limitations}.

\paragraph{«\DEIC{} dynamics: temporal architecture of empathy blocking».} Трёхшкальный временной фрейм (fast/meso/slow), связь с active inference и predictive coding. Эмпирические корреляты через micro-longitudinal (daily diary) и neurocognitive (ERPs) парадигмы.

\paragraph{«ID and P as directional asymmetries of empathic prediction».} Разделение bottom-up ID и top-down P в рамке predictive coding. Interoceptive accuracy tasks как дифференциатор: ID-пациенты должны иметь дефицит интероцептивной точности; P-пациенты --- сохранную интероцепцию при сниженной моральной salience чужого дистресса.

\paragraph{«What does attention do: candidate mechanisms for the A-moderator effect».} Трёхканальная теоретическая работа: temporal resolution, automatic binding decomposition, metacognitive working memory. Empirical probes через phasic attention tasks и dual-task paradigms.

\paragraph{«Individual differences in A-trainability: who benefits from witnessing-based therapy?»} Pre-treatment predictors эффекта A-based interventions. Кандидаты: working memory capacity, interoceptive accuracy, language-mediation, safety-context stability. Клиническая выборка с RCT arm.

\paragraph{«Developmental branching in $2\times 2$ psychopathic topology».} Longitudinal child sample с early CA/temperament-маркерами; тест Karpman-гипотезы о двух путях в первичную vs вторичную клетку.

\paragraph{«\DEIC{} and the reflexivity of measurement: observer-A as part of the instrument».} Методологическая статья о том, что A наблюдателя --- условие возможности видеть A у наблюдаемого. Эпистемологически: структурная, не мистическая аналогия с observer-problem в квантовой механике. Операциональные следствия: supervision-требования для клиницистов, A-assessment исследователя как часть методологического репортинга.

\paragraph{«Free will as a gradient of attention: an operational reading».} Отдельная философско-психологическая статья, развивающая §\ref{subsec:free_will}. Предлагает переформулировку проблемы свободы воли как проблемы измеримой A-инфраструктуры, с проверяемыми предсказаниями: корреляция ощущения свободного выбора с моментальным A (а не trait-self-efficacy), прирост self-determination-маркеров после A-интервенций, структурная функциональная несвобода в low-A-популяциях. Разворачивает этические следствия (перекалибровка ответственности, прощение как эпистемологическая корректность), без отмены различения между low-A-автоматизмом и A-поддержанным P-выбором.

\paragraph{«Pharmacology as DEIC-modification: a stratification-based prediction map».} Статья, формализующая таблицу «класс веществ --- модификация компонента D/P/A --- предсказание по профилю --- proposed clinical test», развёрнутую в §\ref{subsec:pharmacology_nosology}. Ядро --- дифференциальные предсказания: psilocybin большой эффект на D-dominant, малый/обратный на P-dominant; MDMA уникален как E-ceiling$+$D-dissolve$+$A-preserve; ketamine как «диссоциация для растворения диссоциации» с responders$=$hi-D baseline; opioid dependence как safety-deficit-D-signature. Предлагаемая эмпирическая программа: retrospective re-analysis существующих psilocybin/MDMA RCT с DEIC-profile на baseline как модератором ответа.

\paragraph{«The structural under-sizing of modern mental healthcare: what early-formation D requires».} Клинико-институциональная статья, развёртывающая §\ref{subsec:undersizing}. Ядро: argument for a structurally different scale of intervention for personality-formative early D, с обзором методов адекватного масштаба (длительный интенсивный психоанализ, psychedelic-assisted protocols с extended integration, therapeutic communities, multi-year somatic work) и картой «что сейчас недоступно». Не манифест; эмпирическая программа по proof-of-concept: DEIC-стратификация на baseline как предиктор ответа на symptom-level vs structural-level интервенции.

\paragraph{«DEIC и populational-level predictions: attention economy, politics, collective D».} Статья, разворачивающая §\ref{subsec:political}. Предсказания: (а) интенсивность использования engagement-оптимизированных платформ снижает A в 6--12-месячных лонгитюдах; (б) подростковые когорты с ранним algorithm-exposure имеют измеримо иную DEIC-топологию при том же ACE-load; (в) электоральная поддержка P-стилевой демагогии предсказывается general A-уровнем аудитории; (г) DEIC-профили ex-cult и ex-radicalised лиц характеризуются low A с сохранённым подавленным D. Данные для (а)--(г) требуют cross-cohort и cross-cultural дизайнов.

\paragraph{«Evolutionary mismatch and the biological necessity of integration».} Программная статья, развивающая §\ref{subsec:evolutionary}. Тезис: современный взрослый человек --- первый в истории вида, для которого интеграция есть биологическое требование, а не культурная роскошь. Эмпирическая программа: population-level траектории A-индикаторов (distress tolerance, boredom capacity, sustained attention) в связи с urban engagement-density; инцидентность D-dominant расстройств в зависимости от разрыва между ACE-load и доступностью integration-инфраструктуры; 20--40-летние когортные сравнения когорт с и без структурного доступа к integration-систем.

\paragraph{«Nosological reorganization: DEIC-profile as a trans-diagnostic stratifier».} Эмпирическая статья, проверяющая способность DEIC-профиля на baseline захватывать больший процент дисперсии клинического ответа и траектории, чем DSM-метка, в large-scale retrospective re-analysis. Совместимость и параллель с HiTOP и RDoC; falsification: если DEIC-профиль не даёт прироста над diagnosis label --- рамка ограничивается более узкими применениями, нозологическая претензия снимается.

\subsection{Прикладные направления}

\paragraph{Клинические RCT.} Head-to-head тесты: non-dual mindfulness vs MBSR в работе с ID-доминантными клиентами; A-training arm vs 12-step в substance-use с блокирующим профилем; sequenced A-work $\to$ CBT vs CBT-first в treatment-resistant депрессии.

\paragraph{DEIC как решение проблемы под-диагностирования (collaborative-filtering подход).} Прямое расширение §\ref{subsec:dx_count_bias}. Наблюдение, что инструмент систематически ловит именно тех, кого self-report диагностика пропускает (hi-P --- через ROC AUC $< 0.5$ на mood/anxiety/neurodev; hi-ID --- через help-seeking block; ранне-травматический hi-CA --- через «не знаю» как дис­со­циа­тив­ный сигнал), переворачивается в позитивный дизайн. Классическая нозология предсказывает диагноз от сообщённых симптомов; DEIC-coordinates устойчивы в архитектурно-похожих группах, поэтому k-NN / latent-class similarity на 6-факторном профиле позволяет предсказывать \emph{вероятностный клинический ландшафт} для нового респондента по близости к подвыборке с известными диагнозами --- стандартная формулировка collaborative filtering, но с клинической целевой переменной. Дизайн проверки: wave с recruitment через первичку и образовательные институции (вне мем-канала, чтобы разорвать selection bias), DEIC на входе, структурированное клиническое интервью (SCID/MINI) в слепую относительно DEIC-профиля как gold standard; метрики --- sensitivity/specificity DEIC-скрининга к клиническим меткам, incremental AUC над симптомным скринингом (PHQ, GAD), калибрационные кривые для «архитектурно обусловленного» прогноза. Публичный смысл тоньше, чем звучит: цель --- не автоматическая маркировка людей, а снятие bottleneck на распознавание там, где сам механизм патологии блокирует распознавание. Ограничения очевидны (риск false-positive разметки, этика скрининга без согласия, необходимость human-in-the-loop клинициста) и должны быть предметом отдельной прикладной рамки.

\paragraph{Публичная коммуникация.} Отдельная рамка «травма не приговор» разработана вне академического слоя и требует собственной эмпирической валидизации на non-academic аудитории. Тезисы: integrated-survivor как эмпирически-обоснованный горизонт восстановления; различение wellness-industry toxicity от contemplative substance; A-work как путь, а не обещание счастья.

\paragraph{Политическая и организационная коммуникация.} P-style демагогия и low-A аудитории --- testable hypothesis с cross-sectional audience studies. Организационная психология: применимость условной-архитектурной рамки в подборе и supervision ролей, требующих эмоциональной регуляции (медицина, helping professions, силовые структуры).

\paragraph{Образование.} A-oriented школьные программы как primary prevention для CA-высоких когорт. Связь с existing SEL (social-emotional learning) инициативами; эмпирическая проверка, работают ли SEL-программы на механизм, который \DEIC{} постулирует.

\paragraph{Attention-protective институциональная инфраструктура.} Прикладное продолжение §\ref{subsec:political}. Дизайн и оценка interventions уровня дизайна медиа-продуктов (slow-feed, pause-интерфейсы, edit-mode), уровня школьных программ (contemplative education как компонент общего curriculum), уровня регуляции (engagement-оптимизация как public health issue). Testable: дают ли такие изменения измеримый прирост A в популяции при прочих равных.

\paragraph{Свидетельская работа коллективной D: transitional justice как DEIC-операция.} Прикладная статья о том, как DEIC-координаты могут давать количественную меру тому, что transitional justice и memorial practices делают интуитивно: коллективная D-work через создание публичного пространства для аффективной интеграции исторической реальности. Пилотные данные --- на post-tоталитарных обществах с доступными longitudinal social-psychological данными.

\subsection{Мета-программа}

\paragraph{Методологическая статья программы.} Формальное изложение дисциплинарных правил анти-Уилбер (§\ref{subsec:discipline}), reflexivity clause, observer-A problem. Место: статья \#18 в программе, выходит после эмпирического ядра.

\paragraph{Открытая документация программы.} \texttt{PROGRAM\_STRUCTURE.md} в репозитории содержит полный список 18 планируемых статей с их falsifiers и expected deliverables. Каждая замыкается либо pre-registered replication, либо явным обновлением программы при non-replication. Это и есть форма, в которой мы принимаем обязательства --- распределённая во времени сеть проверяемых утверждений, а не завершённая теория.

Текущая статья --- узкая по тону, скромная в претензиях --- защищает возможность развернуть амбициозную программу. \DEIC{} как каркас стоит на том, что каждая из его будущих работ может быть опровергнута в своей точке, и если значительная их часть опровергнется, каркас как целое подлежит пересмотру. Иначе мы работаем не в психометрии, а в риторике.
